\documentclass[a4paper]{article}

% --- Основные пакеты ---
\usepackage{graphicx}     % картинки
\usepackage{geometry}     % поля страницы
\usepackage{hyperref}     % гиперссылки
\usepackage{amsmath}      % математика
\usepackage{amssymb}      % математические символы
\usepackage{float}        % [H] позиционирование фигур
\usepackage[T2A]{fontenc} % кириллица
\usepackage[utf8]{inputenc} % кодировка UTF-8
\usepackage[russian]{babel} % многоязычность
\usepackage{caption}      % настройка подписей
\usepackage{longtable}    % длинные таблицы
\usepackage{array}        % расширенные таблицы
\usepackage{multirow}     % объединение строк в таблицах

% --- Настройка страницы ---
\geometry{top=1.5cm, bottom=1.5cm, left=3cm, right=1.5cm}

\renewcommand{\thesubsection}{\arabic{subsection}.}

% --- Настройка гиперссылок ---
\hypersetup{
    colorlinks,
    linkcolor=black,
    urlcolor=blue,
    citecolor=black
}

% --- Настройка подписей ---
\renewcommand{\figurename}{Рис.}
\captionsetup{labelsep=period}

\begin{document}

% --- Титульная страница ---
\centering
Министерство науки и высшего образования Российской Федерации\\
Федеральное государственное автономное образовательное учреждение\\
высшего образования\\
«Южно-Уральский государственный университет\\
(национальный исследовательский университет)»\\
Институт естественных и точных наук

\setlength{\parskip}{7.5cm}

Отчет\\
Лабораторная работа №7.\\
Решение нелинейных уравнений.

\setlength{\parskip}{6.5cm}

\raggedleft
\begin{minipage}{0.345\textwidth}
\raggedright
Выполнил: студент группы ЕТ-210\\
Тельнов С.М.\\
Проверил: преподаватель кафедры ВМ\\
Стариков Я.Е.
\end{minipage}

\setlength{\parskip}{6cm}

\centering
Челябинск — 2026

\setlength{\parskip}{0pt plus 1pt}
\newpage

% --- Содержание ---
\centering
\textbf{Краткое содержание}
\begin{itemize}
    \setlength{\parskip}{0.4cm}
    \item \hyperref[sec:theory]{\text{Теоретическая часть} \dotfill \pageref{sec:theory}}
    \item \hyperref[sec:practice]{\text{Практическая часть} \dotfill \pageref{sec:practice}}
    \item \hyperref[sec:conclusion]{\text{Вывод} \dotfill \pageref{sec:conclusion}}
\end{itemize}

\setlength{\parindent}{1.5em}

\raggedright

% --- Теоретическая часть ---
\section*{\centering\fontsize{14pt}{0pt}\selectfont\textbf{Теоретическая часть}\label{sec:theory}}

Целью данной работы является ознакомление с итерационными методами решения нелинейных алгебраических и трансцендентных уравнений, а также систем нелинейных уравнений, программная реализация данных алгоритмов и анализ их сходимости.

\subsection{Решение нелинейных уравнений с одной переменной}
Для поиска корней нелинейного уравнения $F(x) = 0$ на заданном отрезке локализации $[a, b]$ применяются следующие приближенные численные методы:
\begin{enumerate}
    \item \textbf{Метод деления отрезка пополам (дихотомии).} На каждом шаге отрезок делится пополам точкой $x = \frac{a+b}{2}$. Из двух полученных половин выбирается та, на концах которой функция меняет знак. Метод абсолютно сходящийся, но обладает невысокой линейной скоростью сходимости.
    \item \textbf{Метод хорд.} Приближение к корню на $k$-й итерации определяется как точка пересечения оси абсцисс прямой (хордой), проведенной через концы текущего интервала:
    \begin{equation}
        x_{k}=a-\frac{b-a}{F(b)-F(a)}F(a)
    \end{equation}
    После чего интервал сужается аналогично методу дихотомии.
    \item \textbf{Комбинированный метод секущих и хорд.} Позволяет ускорить процесс вычислений, используя информацию о поведении функции на предыдущих шагах. Итерационный шаг вычисляется по формуле:
    \begin{equation}
        x_{n+1}=x_{n}-\frac{x_{n}-x_{n-1}}{F(x_{n})-F(x_{n-1})}F(x_{n})
    \end{equation}
    \item \textbf{Метод Хэлли.} Метод третьего порядка сходимости, использующий не только первую, но и вторую производную функции, аппроксимируя её касательной параболой:
    \begin{equation}
        x_{k+1}=x_{k}-\frac{F(x_{k})F'(x_{k})}{(F'(x_{k}))^{2}-0.5F(x_{k})F''(x_{k})}
    \end{equation}
\end{enumerate}

\subsection{Решение систем нелинейных уравнений}
Для решения систем нелинейных уравнений вида $\overline{F}(\overline{X}) = 0$ используются многомерные аналоги итерационных методов.
\begin{enumerate}
    \item \textbf{Метод простой итерации и метод Зейделя.} Система приводится к релаксационному виду $\overline{X}^{(n+1)} = \overline{X}^{(n)} + \mathbf{W}\overline{F}(\overline{X}^{(n)})$, где $\mathbf{W}$ — матрица релаксационных параметров. В методе Зейделя вновь вычисленные компоненты вектора сразу используются на текущем шаге.
    \item \textbf{Метод Ньютона (Ньютона-Рафсона).} Наиболее распространенный метод, обладающий квадратичной скоростью сходимости в окрестности корня. Итерационная процедура имеет вид:
    \begin{equation}
        \overline{X}^{(n+1)} = \overline{X}^{(n)} - \mathbf{J}^{-1}(\overline{X}^{(n)}) \cdot \overline{F}(\overline{X}^{(n)})
    \end{equation}
    где $\mathbf{J}(\overline{X})$ — матрица Якоби (якобиан), элементами которой являются частные производные функций системы по всем независимым переменным. На практике якобиан может вычисляться конечно-разностным способом (численно) или задаваться аналитически.
\end{enumerate}

\subsection{Аналитическое вычисление матрицы Якоби}
Для исследуемой в работе системы уравнений (Вариант 1) аналитический вид матрицы Якоби формируется следующим образом:
\begin{equation}
\begin{cases}
f_1(x_1, x_2) = \sin(x_1 + x_2) - x_2 - 1.2 = 0 \\
f_2(x_1, x_2) = 2x_1 + \cos x_2 - 2 = 0
\end{cases}
\Rightarrow
\end{equation}

\begin{equation}
\Rightarrow
\begin{cases}
J(x_1, x_2) = 
\begin{pmatrix}
\dfrac{\partial f_1}{\partial x_1} & \dfrac{\partial f_1}{\partial x_2} \\
\\[-0.5em]
\dfrac{\partial f_2}{\partial x_1} & \dfrac{\partial f_2}{\partial x_2}
\end{pmatrix}\\
\frac{\partial f_1}{\partial x_1} = \frac{\partial}{\partial x_1} \big(\sin(x_1 + x_2) - x_2 - 1.2\big) = \cos(x_1 + x_2) \\
\frac{\partial f_1}{\partial x_2} = \frac{\partial}{\partial x_2} \big(\sin(x_1 + x_2) - x_2 - 1.2\big) = \cos(x_1 + x_2) - 1 \\
\frac{\partial f_2}{\partial x_1} = \frac{\partial}{\partial x_1} \big(2x_1 + \cos x_2 - 2\big) = 2 \\
\frac{\partial f_2}{\partial x_2} = \frac{\partial}{\partial x_2} \big(2x_1 + \cos x_2 - 2\big) = -\sin x_2
\end{cases}
\Rightarrow
\end{equation}

\begin{equation}
\Rightarrow J(x_1, x_2) = \begin{pmatrix}
\cos(x_1 + x_2) & \cos(x_1 + x_2) - 1 \\
2 & -\sin x_2
\end{pmatrix}
\end{equation}

Критерием остановки для всех итерационных процессов служит достижение заданной точности $\epsilon$, оцениваемой по норме разности векторов приращений или по норме вектора невязки функций $\|\overline{F}(\overline{X})\| < \epsilon$.

\newpage

% --- Практическая часть ---
\section*{\centering\fontsize{14pt}{0pt}\selectfont\textbf{Практическая часть}\label{sec:practice}}

\subsection*{Задача 1. Решение нелинейного уравнения с одной переменной}

Рассматривается трансцендентное уравнение:
\begin{equation}
f_1(x) = \ln x - \tan x = 0
\end{equation}

Из предварительного анализа (графика функции) были определены два интервала локализации корней: $[4.0,\;4.2]$ и $[7.3,\;7.5]$. Для поиска корней на этих интервалах применялись три метода: метод Ньютона, метод деления отрезка пополам (дихотомии) и метод Хэлли. Результаты при точности $\epsilon = 10^{-6}$ представлены в таблице~\ref{tab:task1_1}.

\begin{table}[H]
\centering
\caption{Сравнение методов для $\epsilon = 10^{-6}$}
\label{tab:task1_1}
\begin{tabular}{|c|l|c|c|}
\hline
Интервал & Метод & $x$ & Итераций \\
\hline
\multirow{3}{*}{$[4.0,\;4.2]$}
& Ньютон      & $4.0954617425$ & 3 \\
& Дихотомия   & $4.0954612732$ & 17 \\
& Хэлли       & $4.0954616062$ & 2 \\
\hline
\multirow{3}{*}{$[7.3,\;7.5]$}
& Ньютон      & $7.3903695711$ & 4 \\
& Дихотомия   & $7.3903694153$ & 17 \\
& Хэлли       & $7.3903695719$ & 2 \\
\hline
\end{tabular}
\end{table}

Как видно из таблицы, метод Хэлли демонстрирует наилучшую скорость сходимости (2 итерации на каждом интервале). Метод Ньютона также сходится быстро (3--4 итерации). Метод дихотомии сходится значительно медленнее (17 итераций), что объясняется его линейной скоростью сходимости.

\subsection*{Задача 1.2. Зависимость числа итераций от точности}

Для максимального корня (интервал $[7.3,\;7.5]$) была исследована зависимость числа итераций от заданной точности $\epsilon$. Результаты приведены в таблице~\ref{tab:task1_2}.

\begin{table}[H]
\centering
\caption{Зависимость числа итераций от точности $\epsilon$ (интервал $[7.3,\;7.5]$)}
\label{tab:task1_2}
\begin{tabular}{|c|c|c|c|}
\hline
$\epsilon$ & Ньютон & Дихотомия & Хэлли \\
\hline
$10^{-3}$  & 3  & 8  & 1 \\
$10^{-6}$  & 4  & 17 & 2 \\
$10^{-9}$  & 4  & 27 & 3 \\
$10^{-12}$ & 5  & 37 & 3 \\
$10^{-15}$ & 100 (max) & 100 (max) & 100 (max) \\
\hline
\end{tabular}
\end{table}

Можно наблюдать, что при ужесточении требований к точности число итераций метода дихотомии растёт пропорционально $\log_2(1/\epsilon)$. Методы Ньютона и Хэлли демонстрируют значительно более слабую зависимость, что подтверждает их квадратичную и кубическую сходимость соответственно. При $\epsilon = 10^{-15}$ все методы исчерпали лимит в 100 итераций, что связано с достижением машинной точности представления чисел с плавающей точкой.

\subsection*{Задача 1.3. Метод простой итерации}

Для решения уравнения $f_1(x) = 0$ был применён метод простой итерации с различными значениями релаксационного параметра $\omega$ (матрица $W = [\omega]$). Результаты приведены в таблице~\ref{tab:task1_3}.

\begin{table}[H]
\centering
\caption{Метод простой итерации при $\epsilon = 10^{-6}$}
\label{tab:task1_3}
\begin{tabular}{|c|c|c|c|c|}
\hline
\multirow{2}{*}{$\omega$} & \multicolumn{2}{c|}{$[4.0,\;4.2]$} & \multicolumn{2}{c|}{$[7.3,\;7.5]$} \\
\cline{2-5}
& $x$ & Итераций & $x$ & Итераций \\
\hline
0.1 & 4.0954612930 & 40  & 7.3903694649 & 21 \\
0.2 & 4.0954614322 & 17  & 7.3903695516 & 5  \\
0.3 & 4.0954614241 & 8   & 7.3903693742 & 16 \\
0.4 & 4.0954612627 & 5   & --- & не сошёлся \\
0.5 & 4.0954617466 & 13  & --- & не сошёлся \\
0.6 & 4.0954612664 & 28  & --- & не сошёлся \\
0.7 & --- & не сошёлся & --- & не сошёлся \\
\hline
\end{tabular}
\end{table}

Анализ показывает, что эффективность метода простой итерации существенно зависит от выбора параметра $\omega$. Для интервала $[4.0,\;4.2]$ оптимальным оказалось значение $\omega = 0.4$ (5 итераций), для интервала $[7.3,\;7.5]$ — $\omega = 0.2$ (5 итераций). При $\omega \ge 0.5$ для второго интервала метод расходится, что характерно для методов релаксационного типа при нарушении условий сходимости.

\subsection*{Задача 1.4. Уточнение корней методом fsolve}

Для верификации полученных результатов корни были найдены с помощью встроенной функции \texttt{fsolve} библиотеки SciPy:
\begin{itemize}
    \item Корень в $[4.0,\;4.2]$: $x = 4.095461605909919$;
    \item Корень в $[7.3,\;7.5]$: $x = 7.390369571112944$.
\end{itemize}

Полученные значения хорошо согласуются с результатами, найденными реализованными методами. Следует отметить, что точность \texttt{fsolve} сопоставима с точностью метода Хэлли.

\subsection*{Задача 2. Решение системы нелинейных уравнений}

Исследуется система двух нелинейных уравнений:
\begin{equation}
\begin{cases}
f_1(x_1, x_2) = \sin(x_1 + x_2) - x_2 - 1.2 = 0,\\
f_2(x_1, x_2) = 2x_1 + \cos x_2 - 2 = 0.
\end{cases}
\end{equation}

Область поиска корня: $x_1 \in [1,\;1.3]$, $x_2 \in [-2,\;-1.7]$.

Для решения применялся метод Ньютона с различными значениями точности $\epsilon$. Результаты представлены в таблице~\ref{tab:task2}.

\begin{table}[H]
\centering
\caption{Метод Ньютона для системы нелинейных уравнений}
\label{tab:task2}
\begin{tabular}{|c|c|c|}
\hline
$\epsilon$ & $(x_1,\;x_2)$ & Итераций \\
\hline
$10^{-3}$  & $(1.1407533892,\;-1.8558549954)$ & 2 \\
$10^{-6}$  & $(1.1405671800,\;-1.8557722639)$ & 3 \\
$10^{-9}$  & $(1.1405671770,\;-1.8557722587)$ & 4 \\
$10^{-12}$ & $(1.1405671770,\;-1.8557722587)$ & 4 \\
$10^{-15}$ & $(1.1405671770,\;-1.8557722587)$ & 5 \\
\hline
\end{tabular}
\end{table}

Метод Ньютона демонстрирует быструю сходимость: уже при $\epsilon = 10^{-3}$ достаточно 2 итераций. Начиная с $\epsilon = 10^{-9}$, значения корня стабилизируются на уровне $(1.1405671770,\;-1.8557722587)$, что соответствует машинной точности.

Для верификации решение было найдено с помощью \texttt{fsolve}: $x = (1.1405671770,\;-1.8557722587)$, что полностью совпадает с результатом метода Ньютона.

\subsection*{Задача 2.3. Численная и аналитическая матрица Якоби}

Проведено сравнение метода Ньютона с численным (конечно-разностным) вычислением матрицы Якоби и с её аналитическим заданием. Результаты представлены в таблице~\ref{tab:task2_3}.

\begin{table}[H]
\centering
\caption{Сравнение численного и аналитического якобиана}
\label{tab:task2_3}
\begin{tabular}{|c|c|c|c|}
\hline
\multirow{2}{*}{$\epsilon$} & Численный & \multicolumn{2}{c|}{Аналитический} \\
\cline{2-4}
& Итераций & Итераций & $(x_1,\;x_2)$ \\
\hline
$10^{-3}$  & 2 & 2 & $(1.1407533855,\;-1.8558549952)$ \\
$10^{-6}$  & 3 & 3 & $(1.1405671800,\;-1.8557722639)$ \\
$10^{-9}$  & 4 & 4 & $(1.1405671770,\;-1.8557722587)$ \\
$10^{-12}$ & 4 & 4 & $(1.1405671770,\;-1.8557722587)$ \\
$10^{-15}$ & 5 & 4 & $(1.1405671770,\;-1.8557722587)$ \\
\hline
\end{tabular}
\end{table}

Количество итераций для обоих подходов практически совпадает, что подтверждает корректность численной аппроксимации производных. При $\epsilon = 10^{-15}$ аналитический якобиан даёт на одну итерацию меньше, что связано с отсутствием ошибок численного дифференцирования. Полученные корни идентичны во всех случаях.

\newpage

% --- Вывод ---
\section*{\centering\fontsize{14pt}{0pt}\selectfont\textbf{Вывод}\label{sec:conclusion}}

В ходе выполнения лабораторной работы были реализованы и исследованы численные методы решения нелинейных уравнений и их систем.

\subsection*{Основные результаты}

\begin{enumerate}
    \item \textbf{Решение уравнений с одной переменной.} Для трансцендентного уравнения $f_1(x) = \ln x - \tan x = 0$ найдены два корня на интервалах $[4.0,\;4.2]$ и $[7.3,\;7.5]$. Сравнение методов показало, что метод Хэлли демонстрирует наивысшую скорость сходимости (2 итерации при $\epsilon=10^{-6}$), метод Ньютона — несколько меньшую (3--4 итерации), а метод дихотомии — наименьшую (17 итераций). При ужесточении точности до $\epsilon=10^{-15}$ все методы исчерпали лимит в 100 итераций ввиду ограничений машинной арифметики.
    \item \textbf{Метод простой итерации.} Установлено, что выбор релаксационного параметра $\omega$ критически влияет на сходимость. Оптимальные значения $\omega$ (0.4 для первого интервала и 0.2 для второго) позволили достичь сходимости за 5 итераций. Неудачный выбор $\omega$ приводит к расходимости метода.
    \item \textbf{Системы нелинейных уравнений.} Метод Ньютона успешно решал систему из двух уравнений, демонстрируя быструю сходимость (2--5 итераций в зависимости от $\epsilon$). Результаты верифицированы с помощью встроенной функции \texttt{fsolve}.
    \item \textbf{Численный и аналитический якобиан.} Сравнение показало, что численное (конечно-разностное) вычисление матрицы Якоби даёт практически ту же скорость сходимости, что и аналитическое задание. Лишь при очень высокой точности ($\epsilon=10^{-15}$) аналитический якобиан даёт преимущество в одну итерацию за счёт отсутствия ошибок аппроксимации.
\end{enumerate}

Таким образом, все поставленные задачи выполнены: изучены и программно реализованы основные итерационные методы решения нелинейных уравнений, проведён анализ их сходимости, выполнено сравнение с эталонными решениями. Полученные результаты подтверждают теоретические оценки скорости сходимости рассмотренных методов.

\end{document}